\documentclass[11pt,a4paper]{amsart}
\usepackage[T1]{fontenc}
\usepackage{lmodern}
\usepackage[margin=26mm,headheight=14pt,headsep=7mm]{geometry}
\usepackage{amsmath,amssymb,amsthm,mathtools,mathrsfs}
\usepackage{microtype,etoolbox}
\makeatletter
\patchcmd{\@maketitle}{\global\topskip42\p@}{\global\topskip24\p@}{}{}
\patchcmd{\@maketitle}{\dimen@34\p@}{\dimen@20\p@}{}{}
\makeatother
\usepackage[hidelinks]{hyperref}
\hypersetup{pdftitle={Renyi and Tsallis entropy monotonicity for Bernoulli sums},
 pdfauthor={Kirill Shevtsov},pdfsubject={Hillion--Johnson entropy monotonicity conjecture}}
\numberwithin{equation}{section}
\newtheorem{theorem}{Theorem}[section]
\newtheorem{lemma}[theorem]{Lemma}
\newtheorem{proposition}[theorem]{Proposition}
\newtheorem{corollary}[theorem]{Corollary}
\theoremstyle{remark}
\newtheorem{remark}[theorem]{Remark}
\newcommand{\Ber}{\operatorname{Ber}}
\newcommand{\Bin}{\operatorname{Bin}}
\newcommand{\D}{\mathscr D}
\newcommand{\HH}{\mathscr H}
\newcommand{\eps}{\varepsilon}
\newcommand{\supp}{\operatorname{supp}}
\allowdisplaybreaks[1]
\setlength{\parskip}{1pt}
\emergencystretch=1.5em
\title[Entropy monotonicity for Bernoulli sums]{R\'enyi and Tsallis entropy monotonicity\\for Bernoulli sums}
\author{Kirill Shevtsov}
\address{Faculty of Mathematics, HSE University}
\email{keshevtsov@gmail.com}
\date{November 2, 2026}
\keywords{Poisson--binomial distribution, R\'enyi entropy, Tsallis entropy,
Hillion--Johnson conjecture, log-concavity}
\begin{document}
\begin{abstract}
We prove that the R\'enyi and Tsallis entropies of a sum of independent
Bernoulli random variables are nondecreasing in each success probability
on $[0,1/2]^n$ for every order $0\le q\le2$, as asserted in the
Hillion--Johnson entropy monotonicity conjecture. The proof holds for
arbitrary finite $n$ and is entirely real-variable. A quadratic
interpolation inequality reduces the sign of the entropy derivative
to fair binomial laws. Below order one, the local inequality follows
from two integrating-factor identities for an explicit scalar function;
above order one, it follows from a local third-derivative estimate.
The boundary cases are included, and a two-coin example shows that
monotonicity fails at every order $q>2$. We also obtain a strict
quantitative decrease of Hellinger distance to the uniform law on a
fixed count alphabet under any genuine coordinate increase.
\end{abstract}
\maketitle

\section{Introduction and main result}

For a sum of independent Bernoulli random variables with success
probabilities at most $1/2$, does making one coin closer to fair
increase the entropy of the sum? Hillion and Johnson proved this
for Shannon entropy and formulated the R\'enyi--Tsallis extension
as Conjecture~4.4 in \cite{HJ}. Their paper also treats order zero,
proves the order-two case, and gives counterexamples at every order
greater than two.

We prove the full statement on $0\le q\le2$ for arbitrary finite
Bernoulli sums, also called Poisson--binomial laws. The two open
ranges $0<q<1$ and $1<q<2$ require different local inequalities,
but both lead to the same quadratic interpolation on the parameter cube.

For a finite probability mass function $f=(f_k)$ and $q>0$, $q\ne1$,
the R\'enyi and Tsallis entropies are
\begin{equation}
 H_q(f)=\frac{\log\sum_k f_k^q}{1-q},\qquad
 T_q(f)=\frac{1-\sum_k f_k^q}{q-1}.
 \label{eq:entropies}
\end{equation}
These are the standard families introduced in \cite{Renyi,Tsallis}.
All logarithms are natural. At $q=1$, both are defined as Shannon
entropy,
\[
 H_1(f)=T_1(f)=-\sum_k f_k\log f_k,
 \qquad 0\log0=0.
\]
At order zero, we set $H_0(f)=\log|\supp f|$ and
$T_0(f)=|\supp f|-1$, without any convention for $0^0$.

\begin{theorem}[Entropy monotonicity]\label{thm:main}
Let $n\ge1$. For $p\in[0,1/2]^n$, let $f_p$ be the law of
$\sum_{i=1}^n X_i$, where the $X_i$ are independent and
$X_i\sim\Ber(p_i)$. If
\[
 0\le p_i\le p_i'\le\tfrac12\qquad(1\le i\le n),
\]
then, for every real $0\le q\le2$,
\[
 H_q(f_p)\le H_q(f_{p'}),\qquad
 T_q(f_p)\le T_q(f_{p'}).
\]
The range is sharp: for every real $q>2$, both inequalities fail for
some coordinatewise ordered pair in $[0,1/2]^2$.
\end{theorem}

The proof first reduces coordinate monotonicity to a sign at a fair
moving coin. For a finite log-concave background, we then establish a
quadratic lower bound by its two endpoint values. Iterating this
bound replaces each background coin by either a deterministic zero
or a fair coin. Every terminal law is a fair binomial, and symmetry
makes its fair-coordinate derivative vanish.

Section~\ref{sec:scalar} proves the scalar estimate for orders below
one by elementary differential identities and positive real
integrals. Section~\ref{sec:lowerbellman} turns it into the required
interpolation inequality. Section~\ref{sec:upper} proves the same
interpolation for orders between one and two. The endpoint orders
and sharpness are treated in Section~\ref{sec:scope}, followed by the
Hellinger consequence in Section~\ref{sec:hellinger}.

\section{The fair coordinate reduction}
\label{sec:reduction}

\subsection{It suffices to control a fair coordinate}

Fix all coins but one, and call the remaining background law $g$.
Extend every finite sequence by zero to all integer indices. The
moving law is
\[
f_k(r)=(1-r)g_k+rg_{k-1},\qquad w_k=g_{k-1}-g_k.
\]
For $0<r<1$, with identically zero summands omitted,
\begin{equation}
\frac{d^2}{dr^2}T_q(f(r))
=-q\sum_k w_k^2f_k(r)^{q-2}\le0.
\label{eq:coordinateconcavity}
\end{equation}
The same formula at $q=1$ is $-\sum_k w_k^2/f_k(r)$. Thus the derivative
at $r\le1/2$ is at least its value at $1/2$. Write that fair-coordinate
derivative as
\begin{equation}
\D_q(g)=-\frac{q2^{1-q}}{q-1}
\sum_k(g_{k-1}-g_k)(g_{k-1}+g_k)^{q-1},\qquad q\ne1,
\label{eq:fairD}
\end{equation}
and
\begin{equation}
\D_1(g)=-\sum_k(g_{k-1}-g_k)\log(g_{k-1}+g_k).
\label{eq:fairD1}
\end{equation}
Here \eqref{eq:fairD} is used only for $q>0$, $q\ne1$.
The term at a zero pair is zero. These are continuous functionals on
finite nonnegative sequences: for $q>0$, $q\ne1$, the magnitude of a raw
summand is at most $(g_{k-1}+g_k)^q$, and at $q=1$ use
$x|\log x|\to0$ as $x\downarrow0$.

It is enough to show $\D_q(g)\ge0$ for every lower-cube
Poisson--binomial background. R\'enyi and Tsallis have the same
derivative sign, since
\begin{equation}
\frac{dH_q}{dr}=\frac{1}{\sum_k f_k^q}\frac{dT_q}{dr}.
\label{eq:renyisign}
\end{equation}
For each fixed positive order, entropy is continuous on the fixed
finite alphabet, including zero masses. An interior-coordinate
monotonicity inequality therefore extends to the closed moving
interval by continuity. No finite derivative at moving parameter
zero is required.

\subsection{Log-concavity and the vertices of the parameter cube}

A nonnegative sequence is log-concave if $h_k^2\ge h_{k-1}h_{k+1}$ and its positive
support is an interval. Bernoulli convolution preserves this property.
Indeed, if $g_k=a h_k+b h_{k-1}$ with $a,b\ge0$, then
\begin{align}
g_k^2-g_{k-1}g_{k+1}
={}&a^2(h_k^2-h_{k-1}h_{k+1})
 +b^2(h_{k-1}^2-h_{k-2}h_k)\notag\\
&+ab(h_{k-1}h_k-h_{k-2}h_{k+1})\ge0.
\label{eq:LCconvolution}
\end{align}
The mixed determinant is nonnegative by decreasing consecutive ratios
on the positive support, and directly at the boundary. Starting with a
point mass proves log-concavity for all Bernoulli sums.

At a background parameter vector in $\{0,1/2\}^m$, deleting the zero
coins leaves a fair binomial law. Reflection of its edges preserves
$g_{k-1}+g_k$ and negates $g_{k-1}-g_k$. Therefore
\begin{equation}
\D_q\bigl(\Bin(j,1/2)\bigr)=0\qquad(j\ge0).
\label{eq:vertices}
\end{equation}
We will reduce every background to these vertices. The two ranges of
$q$ use different local inequalities.

\section{A scalar inequality for orders below one}
\label{sec:scalar}

Fix $0<q<1$. For $-1<z<1$ and $0\le s\le1$, define
\begin{align}
 \delta(z)&=\left(\frac{1+z}{2}\right)^q
             -\left(\frac{1-z}{2}\right)^q,
 &h_q(z)&=\frac{z}{\delta(z)},\quad h_q(0)=\frac{2^{q-1}}q,
 \label{eq:delta}\\
 B_s(z)&=(1-sz)^q-1-s^2\bigl[(1-z)^q-1\bigr],
 &K_s(z)&=h_q(z)B_s(z).
 \label{eq:scalarK}
\end{align}
The function $h_q$ is smooth and even. We shall prove the following
shape property on the whole interval, including its negative half.

\begin{lemma}[Scalar shape inequality]\label{lem:scalar}
For $0<q<1$ and $0\le s\le1$,
\begin{equation}
 K_s'''(z)\ge0\quad(-1<z<1),\qquad K_s(0)=K_s''(0)=0.
 \label{eq:Kshape}
\end{equation}
\end{lemma}

The proof uses an auxiliary function whose positivity is obtained
from an integrating factor. Throughout this section, put
\[
 \rho=1-q\in(0,1),\qquad a(x)=1-x^2,\qquad
 c=\rho(2-\rho)=1-q^2.
\]

\begin{lemma}[Auxiliary monotonicity]\label{lem:auxiliary}
For $|x|<1$, let
\begin{equation}
 F(x)=\frac12\int_{-1}^{1}(1+tx)^{-\rho}\,dt
 =\frac{(1+x)^q-(1-x)^q}{2qx},\qquad F(0)=1.
 \label{eq:realF}
\end{equation}
For $0<x<1$, define
\[
 d(x)=\frac{\rho}{a(x)}-\frac{F'(x)}{xF(x)},\qquad
 W(x)=a(x)^\rho F(x)^2,\qquad H(x)=a(x)F(x).
\]
The function $d$ has a smooth even extension at zero, with
$d(0)=c/3$. Moreover,
\[
 d(x)>0,\quad d'(x)>0,\quad H(x)>0,\quad H'(x)<0
 \qquad(0<x<1).
\]
\end{lemma}

\begin{proof}
The integral in \eqref{eq:realF} follows by elementary integration.
It shows that $F$ is positive, even, and smooth on $(-1,1)$.
Indeed, on every compact interval $|x|\le b<1$, the base $1+tx$
is bounded below by $1-b$, so each derivative of the integrand is
uniformly bounded for $-1\le t\le1$. Differentiation under the
integral gives
\[
 F(0)=1,\qquad F'(0)=0,\qquad F''(0)=\frac{\rho(1+\rho)}3.
\]
The identity
\[
 \frac{F'(x)}x=\int_0^1 F''(tx)\,dt
\]
proves smooth removability at zero. Consequently
$d(0)=\rho-F''(0)=c/3$ and $d'(0)=0$.

Consider the differential operator
\[
 L_qy=a y''+2(q-1)xy'-q(q-1)y.
\]
Direct differentiation gives, whenever $1+bx>0$,
\begin{equation}
 L_q[(1+bx)^q]=q(q-1)(b^2-1)(1+bx)^{q-2}.
 \label{eq:poweroperator}
\end{equation}
In particular $L_q[xF]=0$. Expanding this identity yields
\begin{equation}
 aF''+\left[\frac2x-2(1+\rho)x\right]F'
       -\rho(1+\rho)F=0\qquad(x\ne0).
 \label{eq:Fode}
\end{equation}
Set $u=F'/F=x(\rho/a-d)$. Dividing \eqref{eq:Fode} by $F$ and
substituting this expression gives
\[
 0=a(u'+u^2)+[2/x-2(1+\rho)x]u-\rho(1+\rho)
   =\frac ca-3ad-axd'+ax^2d^2.
\]
Thus
\begin{equation}
 xd'+3d=x^2d^2+\frac c{a^2}.
 \label{eq:riccati}
\end{equation}
Differentiating the equality \eqref{eq:riccati}, rather than a sign
inequality, gives
\[
 d''+\left(\frac4x-2xd\right)d'=2d^2+\frac{4c}{a^3}.
\]
The definition of $W$ implies $W'/W=-2xd$, and hence
\begin{equation}
 (x^4Wd')'=x^4W\left(2d^2+\frac{4c}{a^3}\right).
 \label{eq:IFone}
\end{equation}
Since $d'(x)=O(x)$ and $W(0)=1$, the left-hand primitive tends to
zero at the origin. The right-hand side is $O(x^4)$ there. Integrating
\eqref{eq:IFone} from zero therefore gives the exact formula
\begin{equation}
 d'(x)=\frac{\displaystyle\int_0^x t^4W(t)
       \left(2d(t)^2+\frac{4c}{(1-t^2)^3}\right)dt}
       {x^4W(x)}>0.
 \label{eq:positiveD}
\end{equation}
No sign of $d$ was needed to obtain this inequality: the only such
term in the integrand is a square. Since $d(0)=c/3>0$, we also have
$d>0$. Finally,
\begin{equation}
 \frac{H'}H=-x\left(\frac{2-\rho}{a}+d\right)<0,
 \label{eq:Hdecrease}
\end{equation}
while $H>0$ follows directly from $a,F>0$.
\end{proof}

\begin{proof}[Proof of Lemma~\ref{lem:scalar}]
The cases $s=0,1$ give $K_s\equiv0$, so assume $0<s<1$.
By \eqref{eq:realF},
\[
 \widehat K_s(z):=2^{1-q}K_s(z)=\frac{B_s(z)}{qF(z)}.
\]
Because $F$ is even and $B_s(0)=B_s''(0)=0$,
\[
 \widehat K_s(0)=\widehat K_s''(0)=0,
 \qquad \widehat K_s'(0)=-s(1-s).
\]
For $e\in\{-1,1\}$ and $0\le x<1$, set
\[
 Z_e(x)=e\widehat K_s(ex).
\]
The sign convention ensures that
\begin{equation}
 Z_e'''(x)=\widehat K_s'''(ex),\qquad
 Z_e(0)=Z_e''(0)=0.
 \label{eq:signbranches}
\end{equation}
We now derive the same positive curvature equation for both choices
of $e$.

Equation~\eqref{eq:poweroperator}, applied to every term of $B_s$, gives
\begin{equation}
 L_q[B_s(ex)]=q(q-1)(1-s^2)
                [1-(1-esx)^{q-2}].
 \label{eq:Boperator}
\end{equation}
For any smooth $Z$, equations \eqref{eq:Fode} and
$F'/F=x(\rho/a-d)$ give
\begin{equation}
 L_q[FZ]=aF\{Z''-2d(xZ'-Z)\}.
 \label{eq:productoperator}
\end{equation}
For completeness, the coefficients of $Z'$ and $Z$ on the left are
$2aF'+2(q-1)xF=-2aFxd$ and
$aF''+2(q-1)xF'-q(q-1)F=2aFd$, respectively.
Since $FZ_e=eB_s(ex)/q$, it follows that
\begin{equation}
 Z_e''-2d(xZ_e'-Z_e)=f_e,
 \qquad f_e=\frac{\rho(1-s^2)}H N_e,
 \label{eq:curvatureode}
\end{equation}
where
\[
 N_e(x)=e[(1-esx)^{-1-\rho}-1],\qquad
 N_e'(x)=(1+\rho)s(1-esx)^{-2-\rho}\ge0.
\]
All bases are positive. Since $N_e(0)=0$ and $H'\le0$, we have
\begin{equation}
 f_e\ge0,\qquad
 f_e'=\rho(1-s^2)\left(\frac{N_e'}H-\frac{N_eH'}{H^2}\right)\ge0.
 \label{eq:forcingsigns}
\end{equation}

Put $g_e=Z_e''$ and $V_e=xZ_e'-Z_e$. By \eqref{eq:signbranches},
\[
 V_e(0)=0,\qquad V_e'=xg_e,\qquad g_e=f_e+2dV_e.
\]
Since $W'/W=-2xd$, these identities imply
\[
 (WV_e)'=xWf_e.
\]
Integration, with $W(0)=1$, gives a positive real integral:
\begin{equation}
 V_e(x)=\frac1{W(x)}\int_0^x tW(t)f_e(t)\,dt\ge0.
 \label{eq:volterra}
\end{equation}
Differentiating the exact identity $g_e=f_e+2dV_e$ now yields
\begin{equation}
 g_e'=f_e'+2d'V_e+2dV_e'
      =f_e'+2xd f_e+(2d'+4xd^2)V_e\ge0.
 \label{eq:thirdpositive}
\end{equation}
Each term has a separately established nonnegative sign, by
Lemma~\ref{lem:auxiliary}, \eqref{eq:forcingsigns}, and
\eqref{eq:volterra}. Thus \eqref{eq:signbranches} proves the third
derivative inequality on both halves of $(-1,1)$. Smoothness covers
the origin. Multiplication by $2^{q-1}>0$ returns to $K_s$ and proves
the lemma.
\end{proof}

All integrations above take place on compact subintervals of
$[0,1)$. No limit of the auxiliary differential equations at $x=1$
is needed. The values at $z=\pm1$ required below will instead be
handled by continuity of homogeneous expressions.

\section{Quadratic interpolation for orders below one}
\label{sec:lowerbellman}

For positive $a,b$, define the symmetric, degree-zero function
\begin{equation}
\HH_q(a,b)=h_q\!\left(\frac{a-b}{a+b}\right)
 =\frac{(a-b)(a+b)^{q-1}}{a^q-b^q},
\qquad \HH_q(a,a)=\frac{2^{q-1}}q.
\label{eq:edgeH}
\end{equation}
For a positive log-concave triple $(v,u,t)$, let
\begin{equation}
A=pv+(1-p)u,\quad C=pu+(1-p)t,\quad x=A+C,\quad w=A-C,
\label{eq:localAC}
\end{equation}
and put
\begin{equation}
J(p)=-\HH_q(v,u)A^q+\HH_q(u,t)C^q+w x^{q-1}.
\label{eq:localJ}
\end{equation}

\begin{proposition}[Local quadratic interpolation]
\label{prop:localchord}
If $0<q<1$, $u^2\ge vt$, and $0\le p\le1/2$, then
\begin{equation}
J(p)\ge(1-2p)^2J(0)+4p(1-p)J(1/2).
\label{eq:localchord}
\end{equation}
Also $J(0)=J(1)=u^q[\HH_q(u,t)-\HH_q(v,u)]$.
\end{proposition}

\begin{proof}
The endpoint identity follows from \eqref{eq:edgeH}. For the inequality,
normalize by
\[
L=v+2u+t,\quad z=\frac{v-t}{L},\quad
\kappa=\frac{u^2-vt}{L^2},\quad
\theta=\frac{1+z}{2},\quad\eta=\frac{1-z}{2},\quad s=1-2p.
\]
The two edge scores are
\[
\alpha=z-\frac{2\kappa}{\theta},\qquad
\gamma=z+\frac{2\kappa}{\eta}.
\]
Indeed,
\[
\frac vL=\theta^2-\kappa,\quad
\frac uL=\theta\eta+\kappa,\quad
\frac tL=\eta^2-\kappa,
\]
and $\alpha=(v-u)/(v+u)$, $\gamma=(u-t)/(u+t)$. In particular the
entire interval from $0$ to the given $\kappa$ is feasible at fixed
$L,z$. Moreover,
\begin{align*}
A&=\frac L2\theta(1-s\alpha),&
C&=\frac L2\eta(1-s\gamma),\\
x&=\frac L2(1-sz),&
w&=\frac L2[z(1-sz)+4s\kappa].
\end{align*}
Direct substitution shows that the gap in \eqref{eq:localchord},
divided by $(L/2)^q$, is
\begin{align}
G(\kappa)={}&-\theta^qK_s(\alpha)+\eta^qK_s(\gamma)+zB_s(z)
\notag\\
&+4\kappa\bigl[s(1-sz)^{q-1}-s^2(1-z)^{q-1}\bigr].
\label{eq:normalizedgap}
\end{align}
At $\kappa=0$, both scores are $z$, and $h_q(z)(\theta^q-\eta^q)=z$,
so $G(0)=0$. Differentiating at fixed $z$ gives
\begin{equation}
G'(\kappa)=2\bigl[\theta^{q-1}K_s'(\alpha)
 +\eta^{q-1}K_s'(\gamma)\bigr]-\frac4qB_s'(z).
\label{eq:curvaturederivative}
\end{equation}

Let $S=\theta^{q-1}+\eta^{q-1}$. The weighted mean of the scores is
\[
m=z+\frac{2\kappa}{S}(\eta^{q-2}-\theta^{q-2}).
\]
It moves away from zero: $m\ge z\ge0$ if $z\ge0$, and $m\le z\le0$
if $z\le0$. Lemma~\ref{lem:scalar} says $K_s'$ is convex, and
$K_s''(0)=0$ implies $K_s'$ increases on the positive half and
decreases on the negative half. Jensen's inequality therefore gives
\begin{equation}
\theta^{q-1}K_s'(\alpha)+\eta^{q-1}K_s'(\gamma)
 \ge S K_s'(m)\ge S K_s'(z).
\label{eq:jensen}
\end{equation}

To determine the sign of the resulting lower bound, note that
$\delta'=qS/2$ and $B_s=\delta K_s/z$. Consequently
\begin{equation}
2S K_s'(z)-\frac4qB_s'(z)
=\frac4q[z\delta'(z)-\delta(z)]
  \left(\frac{K_s(z)}z\right)'.
\label{eq:geometricidentity}
\end{equation}
No power series is needed here: direct differentiation gives
\[
\delta''(z)=2^{-q}q(q-1)
 [(1+z)^{q-2}-(1-z)^{q-2}],
\]
which has the sign of $z$. Thus $M(z)=z\delta'(z)-\delta(z)$
satisfies $M(0)=0$ and $M'(z)=z\delta''(z)\ge0$, so $M(z)$
also has the sign of $z$.
Meanwhile $zK_s'-K_s$ vanishes at zero and has derivative
$zK_s''(z)\ge0$ on both sides of zero. It therefore also has the
sign of $z$, as does $(K_s/z)'$. The product in
\eqref{eq:geometricidentity} is nonnegative. At zero the quotient is
smoothly removable, since
\[
 \frac{K_s(z)}z=\int_0^1K_s'(tz)\,dt,\qquad
 \left.\left(\frac{K_s(z)}z\right)'\right|_{z=0}
 =\frac12 K_s''(0)=0.
\]

Equations \eqref{eq:curvaturederivative}--\eqref{eq:geometricidentity}
show $G'(\kappa)\ge0$. Integrating from $0$ to $\kappa$ proves the
proposition.
\end{proof}

\subsection{Zeros and exact telescoping}

The function $h_q$ is bounded and continuous on $[-1,1]$, with
$h_q(\pm1)=1$. Hence
\[
\HH_q(a,b)[pa+(1-p)b]^q
\]
extends continuously to all nonnegative pairs, with value zero at
$(0,0)$. The raw term $w x^{q-1}$ extends as well, since its magnitude
is at most $x^q$. Every nonnegative log-concave triple can be
approximated by positive ones: take $v_\eps=v+\eps$, $t_\eps=t+\eps$,
and $u_\eps=\max\{u,\sqrt{v_\eps t_\eps}\}+\eps$.
Thus Proposition~\ref{prop:localchord} includes zero entries by
continuity.

Now let $h$ be any finite log-concave sequence and
$g_k(p)=(1-p)h_k+ph_{k-1}$. At the triple
$(v,u,t)=(h_{k-2},h_{k-1},h_k)$, the masses in \eqref{eq:localAC} are
$A=g_{k-1}(p)$ and $C=g_k(p)$. The first two terms of
\eqref{eq:localJ} cancel exactly when summed over all $k$.
All terms outside the finite support are zero. Multiplying the summed
inequality by $q2^{1-q}/(1-q)>0$ proves the following.

\begin{theorem}[Below-one Bellman inequality]
\label{thm:lowerbellman}
For $0<q<1$, a finite log-concave $h$, and $0\le p\le1/2$,
\begin{equation}
\D_q(h*\Ber(p))\ge
 (1-2p)^2\D_q(h)+4p(1-p)\D_q(h*\Ber(1/2)).
\label{eq:globalchord}
\end{equation}
\end{theorem}

Both coefficients in \eqref{eq:globalchord} are nonnegative and sum
to one. For a Poisson--binomial background with parameters
$p_1,\ldots,p_m\in[0,1/2]$, apply the inequality successively to each
parameter. The remaining background at each step is log-concave by
\eqref{eq:LCconvolution}. The resulting lower bound is a convex combination of the
values at all vertices $\eps\in\{0,1\}^m$, with weights
\[
W_\eps=\prod_{\eps_i=0}(1-2p_i)^2
        \prod_{\eps_i=1}4p_i(1-p_i).
\]
Each vertex value is zero by \eqref{eq:vertices}, so $\D_q(g)\ge0$.
Equations \eqref{eq:coordinateconcavity} and \eqref{eq:renyisign}
prove Theorem~\ref{thm:main} for $0<q<1$.

\section{Quadratic interpolation for orders between one and two}
\label{sec:upper}

We give a separate proof for $1<q<2$. Fix $0<p<1$ and a positive
log-concave triple $(v,u,t)$. Use $A,C,x,w$ from
\eqref{eq:localAC}, and set
\[
B=v-t,\qquad E=v-2u+t,\qquad T=u^2-vt\ge0.
\]
Define
\begin{align}
\mathcal K_{p,q}(v,u,t)
 &=3EB^2x^{q-3}+(q-3)wB^3x^{q-4},\notag\\
\mathcal H_{p,q}(a,b)
 &=2^{q-1}\frac{(a-b)^3}{[pa+(1-p)b]^{3-q}},\notag\\
R_q&=\mathcal K_{p,q}(v,u,t)
       -\mathcal H_{p,q}(v,u)+\mathcal H_{p,q}(u,t).
\label{eq:upperR}
\end{align}

\begin{lemma}[Upper-range local inequality]
\label{lem:upperlocal}
$R_q\ge0$ for every positive log-concave triple and $1\le q\le2$.
The inequality extends to nonnegative log-concave triples by continuity.
\end{lemma}

\begin{proof}
The elementary identities
\begin{equation}
wB-Ex=2T,\qquad (v-u)x=AB-T,\qquad(u-t)x=CB+T
\label{eq:upperidentities}
\end{equation}
give two exact endpoint formulas:
\begin{equation}
R_1=\frac{T^2\{3wB\,AC+(A^2+C^2)T\}}{x^3A^2C^2},
\qquad R_2=\frac{2T^3}{x^2AC}.
\label{eq:upperendpoints}
\end{equation}
They follow by substituting \eqref{eq:upperidentities} into the two
cubic correction terms and expanding. For example, at $q=1$,
multiply by $x^3A^2C^2$: the resulting identity is
\[
3EB^2xA^2C^2-2wB^3A^2C^2
 -(AB-T)^3C^2+(CB+T)^3A^2
=T^2[3wBAC+(A^2+C^2)T],
\]
using $Ex=wB-2T$, $x=A+C$ and $w=A-C$. At $q=2$,
multiplication by $x^2AC$ similarly gives
\[
3EB^2xAC-wB^3AC
 -2(AB-T)^3C/x+2(CB+T)^3A/x=2T^3.
\]


The formula for $R_2$ is nonnegative. The formula for $R_1$ is immediate
when $wB\ge0$. Otherwise the reflection
$(v,u,t,p)\mapsto(t,u,v,1-p)$ lets us suppose $B>0$ and $w<0$.
Log-concavity and
$w=p(v-u)+(1-p)(u-t)$ then force $u>v>t$. Put $a=u-v$, $b=u-t$.
Here $0<a<b\le u$, $B=b-a$, and $-w\le a$. Hence
\[
T=u(a+b)-ab\ge b^2\ge4a(b-a)\ge4(-wB).
\]
The first difference is $(u-b)(a+b)\ge0$ and the second inequality is
$(b-2a)^2\ge0$. Writing $d=-wB>0$, the bracket in $R_1$ is at least
\[
d[4(A^2+C^2)-3AC]\ge5dAC>0.
\]
This proves the $q=1$ endpoint.

For $1<q<2$, first suppose $E\le0$. For differentiable $\psi$, define
the linear functional
\begin{equation}
\mathcal L(\psi)=3EB^2\psi(x)+wB^3\psi'(x)
-4[(v-u)^3\psi(2A)-(u-t)^3\psi(2C)].
\label{eq:Lfunctional}
\end{equation}
It equals $R_q$ when $\psi(y)=y^{q-3}$. The positive integral
representation
\begin{equation}
y^{q-3}=\frac1{I_q}\int_0^\infty
 r^{q-2}(y+r)^{-2}\,dr,
\quad I_q=\int_0^\infty\frac{r^{q-2}}{(1+r)^2}\,dr>0
\label{eq:shiftrep}
\end{equation}
shows that $R_q$ is a positive weighted integral of $R_1$ evaluated
on the common-shifted triple $(v+r/2,u+r/2,t+r/2)$. Indeed, the
differences $B,E,w$ remain fixed, while $x,2A,2C$ all increase by $r$.
The shifted determinant is $T-rE/2\ge0$, so every such $R_1$ is
nonnegative. Differentiation and the integral interchange are
justified locally uniformly for positive arguments by integrable
dominating powers: near zero the weight is $r^{q-2}$, and at infinity
the integrands are $O(r^{q-4})$ or faster; differentiating a shifted
inverse square only improves its decay.
Thus $R_q\ge0$ if $E\le0$.

The remaining case can be covered whenever $wB\ge0$. Normalize by
\[
z=w/x,\qquad r=B/x,\qquad \tau=T/x^2.
\]
Reflection lets us assume $z,r\ge0$; $|z|<1$ and $\tau\ge0$.
Define
\begin{align*}
U(z)&=2-(1+z)^{q-1}-(1-z)^{q-1},\\
V(z)&=qz-\tfrac12[(1+z)^q-(1-z)^q],\\
W(z)&=(1+z)^{q-3}+(1-z)^{q-3}.
\end{align*}
Expanding the two cubes using \eqref{eq:upperidentities} gives exactly
\begin{equation}
\frac{R_q}{x^q}=V(z)r^3-3U(z)r^2\tau
 +\frac{6U'(z)}{q-1}r\tau^2+4W(z)\tau^3.
\label{eq:uppernormal}
\end{equation}
For $0<z<1$, all four functions $U,U',V,W$ are positive, and
\[
U''=(q-1)(2-q)W>0,\qquad
V(z)=\frac q2\int_0^zU(a)\,da.
\]
Since $U'$ is increasing,
\begin{equation}
U(z)^2=2\int_0^zU(a)U'(a)\,da
 \le\frac4q U'(z)V(z).
\label{eq:UVbound}
\end{equation}
For $r>0$, the first three terms in \eqref{eq:uppernormal} form a
quadratic in $\tau$ with positive constant and quadratic
coefficients. Its squared linear coefficient divided by four times
their product is at most
\[
\frac{9U^2r^4}{4Vr^3\,[6U'r/(q-1)]}
 \le\frac{3(q-1)}{2q}<1.
\]
Thus the quadratic is nonnegative, and so is the last term in
\eqref{eq:uppernormal}. If $z=0$ or $r=0$, the first three terms
vanish and the conclusion is immediate.

These cases exhaust all triples: the identity $Ex=wB-2T$ shows
that $wB<0$ implies $E<0$. This proves $R_q\ge0$.
For fixed $p\in(0,1)$, both the two-point correction and the local
kernel have magnitude bounded by a constant times the $q$th power
of the sum of their arguments. They therefore extend continuously
at zero. Approximation by positive log-concave triples, as in
Section~\ref{sec:lowerbellman}, finishes the proof.
\end{proof}

\subsection{The same quadratic Bellman inequality above one}

Let $h$ be a finite log-concave background and
$g_k(p)=(1-p)h_k+ph_{k-1}$. Write $D(p)=\D_q(g(p))$.
For the local triple at index $k$, $x'=B$ and $w'=E$.
Differentiating \eqref{eq:fairD} three times yields
\begin{equation}
D'''(p)=q(2-q)2^{1-q}
\sum_k\mathcal K_{p,q}(h_{k-2},h_{k-1},h_k)\ge0.
\label{eq:Dthird}
\end{equation}
The inequality follows from Lemma~\ref{lem:upperlocal}, because all
two-point correction terms telescope. The function $D$ is continuous
on $[0,1]$ and $C^3$ on $(0,1)$: every non-identically-zero
adjacent mass sum is positive there. Also $D(0)=D(1)$ by shift
invariance. No derivative regularity at the endpoints is needed.

\begin{lemma}[Three-node quadratic interpolation]
\label{lem:quadraticbridge}
If $D\in C[0,1]\cap C^3(0,1)$, $D'''\ge0$, and $D(0)=D(1)$, then
\[
D(p)\ge (1-2p)^2D(0)+4p(1-p)D(1/2),\qquad0\le p\le1/2.
\]
\end{lemma}
\begin{proof}
Let $Q(p)=(1-2p)^2D(0)+4p(1-p)D(1/2)$, the quadratic
interpolant at $0,1/2,1$. For $0<p<1/2$, subtract from $D-Q$ the
cubic $\lambda t(t-1/2)(t-1)$ chosen to make the difference zero
at $t=p$. This difference has four distinct zeros in $[0,1]$.
Three applications of Rolle's theorem give
$6\lambda=D'''(\xi)\ge0$ for some $\xi\in(0,1)$.
Since $p(p-1/2)(p-1)>0$, $D(p)-Q(p)\ge0$.
At $p=0,1/2$ there is equality by definition.
\end{proof}

Consequently the below-one Bellman inequality \eqref{eq:globalchord}
holds also for every $1<q<2$, with exactly the same coefficients,
for every finite log-concave background. The same product-weight
rounding argument already given after Theorem~\ref{thm:lowerbellman}
therefore applies without modification: iterating over the finite
background parameters gives a convex combination of fair-binomial
vertex values, all zero. Thus $\D_q(g)\ge0$, proving
Theorem~\ref{thm:main} for $1<q<2$ by the common reduction.

\section{Endpoint orders and sharpness}
\label{sec:scope}

At $q=0$, let $r$ be the number of strictly positive Bernoulli
parameters. On the lower cube the sum has support $\{0,\ldots,r\}$,
so $T_0=r$ and $H_0=\log(r+1)$. Increasing a coordinate cannot decrease
$r$, proving this endpoint directly.

At $q=1$, take the pointwise limit $T_q(f)\to-\sum_k f_k\log f_k$ on
each fixed finite support. Passing to the limit in the already proved
monotonicity inequality gives Shannon monotonicity, including all
zero-probability boundary cases.

At $q=2$, the fair derivative is identically zero for every finite
background:
\[
\D_2(g)=-\sum_k(g_{k-1}^2-g_k^2)=0.
\]
Coordinate concavity again gives monotonicity for the moving parameter
below one half. This completes the positive part of
Theorem~\ref{thm:main}.

\subsection{A counterexample for every \texorpdfstring{$q>2$}{q greater than two}}

Fix $q>2$ and set
\[
r=\frac{q-2}{q(q-1)}\in(0,1/2).
\]
The upper bound follows from
$q(q-1)-2(q-2)=q^2-3q+4=(q-3/2)^2+7/4>0$.
For the one-coin background $g=(1-r,r)$, formula \eqref{eq:fairD}
becomes
\[
\D_q(g)=-\frac{q2^{1-q}}{q-1}
 [1-2r-(1-r)^q+r^q].
\]
The third derivative of $(1-r)^q$ in $r$ is negative. Taylor's theorem
therefore gives
\[
(1-r)^q\le1-qr+\frac{q(q-1)}2r^2.
\]
The bracket is consequently at least
\[
(q-2)r-\frac{q(q-1)}2r^2+r^q
=\frac{q-2}{2}r+r^q>0.
\]
Thus $\D_q(g)<0$. Continuity places a negative derivative at moving
parameters strictly below and sufficiently close to $1/2$ as well.
R\'enyi has the same sign by \eqref{eq:renyisign}. For example,
$q=3$ and $r=1/6$ give exactly $\D_3(g)=-5/144$.

\section{A quantitative Hellinger consequence}
\label{sec:hellinger}

The order $q=1/2$ gives a quantitative comparison in Hellinger distance.
The reference distribution is uniform on a fixed alphabet, even when
some Bernoulli parameters vanish.

\begin{corollary}[Strict quantitative improvement toward fixed-alphabet uniform]
\label{cor:hellinger}
Let $n\ge1$, let $f$ be the law of a sum of independent Bernoulli
variables with parameters in $[0,1/2]^n$, and let $U_n$ be uniform on
$\{0,\ldots,n\}$. Use the convention
\[
h^2(f,U_n)=1-\sum_{k=0}^n\sqrt{\frac{f_k}{n+1}}
 =\frac12\sum_{k=0}^n
   \left(\sqrt{f_k}-\frac1{\sqrt{n+1}}\right)^2.
\]
Every genuine coordinate increase strictly decreases this distance.
More precisely, freeze the other $n-1$ coins with law $g$, extend it
by zero to all integers, put $w_k=g_{k-1}-g_k$, and write
$f_k(r)=(1-r)g_k+rg_{k-1}$. For $0\le a<b\le1/2$,
\begin{equation}
h^2(f(a),U_n)-h^2(f(b),U_n)
\ge\frac{(1/2-a)^2-(1/2-b)^2}{8\sqrt{n+1}}
      \sum_k w_k^2>0.
\label{eq:hellingerbound}
\end{equation}
Consequently a coordinatewise ordered pair of distinct parameter
vectors gives strict decrease. Independent coordinate thinning
$p_i\mapsto t_ip_i$, $0\le t_i\le1$, cannot decrease the distance,
and increases it strictly if at least one parameter changes.
\end{corollary}

\begin{proof}
At order $1/2$,
\[
T_{1/2}(f)=2\left(\sum_k\sqrt{f_k}-1\right),
\qquad
h^2(f(a),U_n)-h^2(f(b),U_n)
=\frac{T_{1/2}(f(b))-T_{1/2}(f(a))}{2\sqrt{n+1}}.
\]
For $0<r<1/2$, every non-identically-zero $f_k(r)$ is positive and
at most one. Write $S_g=\sum_k w_k^2$. The exact second derivative
\eqref{eq:coordinateconcavity} gives
\[
\frac{d^2}{dr^2}T_{1/2}(f(r))
=-\frac12\sum_k w_k^2 f_k(r)^{-3/2}\le-\frac12 S_g.
\]
Theorem~\ref{thm:main} at the open order $1/2$, through its established
fair-derivative sign, gives
\[
\left.\frac d{dr}T_{1/2}(f(r))\right|_{r=1/2}
=\D_{1/2}(g)\ge0.
\]
Backward integration yields
\[
\frac d{dr}T_{1/2}(f(r))\ge\frac12(1/2-r)S_g.
\]
Integrating from $a$ to $b$ yields
\[
T_{1/2}(f(b))-T_{1/2}(f(a))
\ge\frac14[(1/2-a)^2-(1/2-b)^2]S_g.
\]
Division by $2\sqrt{n+1}$ proves \eqref{eq:hellingerbound} for $a>0$.
If $a=0$, first integrate from $\eps>0$ and pass to $\eps\downarrow0$
by finite-alphabet continuity; no finite derivative at zero is used.
A finite nonzero probability sequence cannot satisfy
$g_{k-1}=g_k$ for every integer, so $S_g>0$. The square difference is
positive for $a<b\le1/2$, also when $b=1/2$. A finite path of
coordinate changes proves the vector comparison and its thinning
consequence. The reference $U_n$ never changes along this path.
\end{proof}

\begin{thebibliography}{9}
\bibitem{HJ}
E. Hillion and O. Johnson,
\emph{A proof of the Shepp--Olkin entropy monotonicity conjecture},
Electronic Journal of Probability \textbf{24} (2019), paper no.~126, 1--14.
\href{https://doi.org/10.1214/19-EJP380}{doi:10.1214/19-EJP380}.
\href{https://arxiv.org/abs/1810.09791}{arXiv:1810.09791}.
\bibitem{Renyi}
A. R\'enyi,
\emph{On measures of entropy and information},
Proceedings of the Fourth Berkeley Symposium on Mathematical
Statistics and Probability, Vol.~1, University of California Press,
Berkeley, 1961, pp.~547--561.
\bibitem{Tsallis}
C. Tsallis,
\emph{Possible generalization of Boltzmann--Gibbs statistics},
Journal of Statistical Physics \textbf{52} (1988), 479--487.
\href{https://doi.org/10.1007/BF01016429}{doi:10.1007/BF01016429}.
\end{thebibliography}
\end{document}
